La estimación de la edad cerebral a partir de datos de neuroimagen permite cuantificar el envejecimiento del cerebro y detectar desviaciones asociadas a enfermedades neurodegenerativas. La diferencia entre la edad predicha y la edad cronológica ---la \textit{brain age gap}--- constituye un potencial biomarcador de envejecimiento acelerado.

En este trabajo se propone un pipeline multimodal que integra conectividad funcional derivada de fMRI en reposo, comprimida mediante un $\beta$-VAE, con medidas volumétricas de materia gris (T1w) para predecir la edad cerebral con XGBoost. El pipeline fue evaluado sobre $N = 1\,245$ sujetos de la cohorte multicéntrica latinoamericana RedLaT, alcanzando un error absoluto medio (MAE) de 5,62 años. La integración de ambas modalidades superó el uso de cualquiera por separado, y la reducción de dimensionalidad resultó más determinante que la elección entre métodos lineales y no lineales o entre familias de regresores.

Se analizó además el impacto de variables demográficas ---años de educación y sitio de adquisición--- sobre la predicción y la brain age gap. La inclusión de estas variables como features adicionales redujo el MAE a 5,17 años, mientras que el sitio de adquisición emergió como el principal confundidor entre centros. La evaluación del pipeline como biomarcador mediante entrenamiento exclusivo en controles sanos mostró señales prometedoras al incorporar información demográfica, aunque el tamaño de la muestra disponible limitó su aplicabilidad clínica. El pipeline es modular y reutilizable, lo que facilita su adaptación a cohortes más amplias y la incorporación de nuevas variables.

\vspace{1em}
\noindent\textbf{Palabras clave:} edad cerebral, neuroimagen, conectividad funcional, fMRI, autoencoder variacional, XGBoost, datos multimodales, brain age gap, variables demográficas, cohorte latinoamericana.
